\chapter{Introduction}

\section{Problem Statement}

A recurring problem in competitive online games is that players form strong opinions about each other's performance based on partial, subjective, and short-term impressions. After a lost match, blame-shifting arrives quickly, but it is far rarer for players to actually examine the match data, understand what happened, how they themselves performed, and where that performance fits within a broader social network.

\emph{League of Legends} is a particularly interesting setting for this, because it simultaneously offers extremely rich match data, recurring teammate and opponent relationships, and social patterns that can be represented and analysed as a graph. When a player repeatedly ends up in the same game as the same people --- as allies or as opponents --- the resulting relationship network goes beyond the momentary statistics of any single match.

This observation is the starting point of the thesis. The goal: build a player network from real match history, analyse that network using a graph-theoretic perspective, and display the results on a visually interpretable frontend.

\section{Motivation}

The project has partly personal origins. After ranked matches, a recurring question was whether repeated negative encounters were coincidental or had some underlying structure --- perhaps the same people being placed together again and again, in a way the matchmaking interface does not surface. That alone is not a scientific question, but it was enough as a starting point. If the Riot API provides detailed match data and those matches contain recurring player relationships, then a graph can in principle be constructed from them, and where there is a graph, there is something to analyse.

The project scope expanded gradually during development. What began as quick data visualisation grew into a dedicated data-collection layer, a normalised scoring pipeline, SQLite-based persistence, a Rust runtime, and a multi-page React frontend. The Python layer remained for data collection and compatibility. Every extension resolved a specific technical constraint rather than following a pre-planned roadmap.

At one point the graph-analysis pipeline began producing cluster-profile exports, and the question arose whether those profiles could be used as inputs to an agent-based simulation. That is how Tribal NeuroSim emerged: the clusters from both datasets compete under identical simulation conditions.

\section{Objectives}

The thesis documents the design, implementation, and evaluation of a multi-component system built on real match data. The system constructs a player-relationship graph from League of Legends match history: repeated co-play and head-to-head appearances produce edges, to which role-aware performance metrics are attached. Graph analysis covers clustering, core-periphery structure, assortativity, and Brandes betweenness centrality, all in a Rust-based, deterministic, and reproducible implementation.

The player network is examined from three perspectives. From a network perspective: the \textit{flexset} dataset exhibits a core-periphery structure. From a metric perspective: the role-aware \textit{opscore} and \textit{feedscore} measures are computed from real match data. From a systems perspective: each component in the Python, SQLite, Node/Express, Rust, and React stack serves a specific purpose that the other layers cannot fulfil.

Every step in the pipeline is re-runnable with concrete inputs and outputs: from data collection through graph construction and clustering to simulation. The parallel analysis of the Flex Queue and SoloQ datasets makes this measurable --- the same pipeline on two different inputs produces different results, and the difference becomes interpretable through the queue mechanics.

\section{Research Questions}

This thesis addresses the following research questions:

\begin{itemize}
    \item How can a real League of Legends match history be used to build a player network that captures not only co-play but also opponent relationships?
    \item How can player performance be scored in a way that accounts for the actual role played while remaining interpretable?
    \item To what extent is it justified to split the system into Python, SQLite, Node.js, Rust, and React components?
    \item Does the current graph show positive assortativity with respect to \textit{opscore} and \textit{feedscore}?
    \item In what ways do the relationship structures of \textit{flexset} and \textit{soloq} differ, and to what extent can the \textit{flexset} graph be described as an associative core-periphery player graph?
\end{itemize}

\section{Scope and Limitations}

The project is a relatively large system, and it would be easy to present every existing idea as a finished result. The thesis tries to avoid this. What actually runs and can be verified appears as complete. What development plans list as future goals does not appear as a finished feature.

\section{Structure of the Thesis}

The chapters proceed from data collection toward graph analytics and simulation experiments.

The opening chapters establish the starting point: the problem statement, the literature background, and the Riot API-based data sources. This is followed by the system overview, then the \textit{opscore} performance metric, the graph construction and clustering mechanics, the Rust runtime and analytics modules. The comparison of the two datasets (\textit{flexset} and \textit{soloq}) is the most concrete research result of the system. The final chapters cover the database sub-project and the Tribal NeuroSim simulation. This document is the condensed, lite version; the extended version \cite{premadegraph-full} develops in separate chapters the topics only briefly touched on here.

\begin{figure}[H]
    \centering
    \includegraphics[width=0.92\textwidth]{match-analysis-page.png}
    \caption{One entry point of the current frontend: the match analysis and data collection interface}
\end{figure}
